\section{主体性：AI 时代“做自己”的工程化含义}

访谈中最强的概念之一，是“构建主体性”。Tristan 说，AI 时代唯一重要的事情可能是“做自己”，但这里的“做自己”不是心灵鸡汤，而是让 AI 知道你是谁：你的价值观、审美、历史资产、表达方式、当前目标和边界。主体性越清晰，分身越能代表你行动。

\begin{figure}[H]
\centering
\includegraphics[width=0.94\textwidth]{subjectivity-stack.png}
\caption{主体性构建：让 AI 能代表你行动的上下文栈。}
\end{figure}

\begin{knowledgebox}{读图：主体性为什么是工程问题}
图中主体性由历史资产、价值观审美、当前目标、行为边界和 agentic 发生组成。它不是抽象个性，而是 AI 分身能否像你、能否替你做决定、能否知道何时停下的一组输入。
\end{knowledgebox}

\subsection{为什么“做自己”会变重要}

当 AI 能主动写文章、发评论、识别人、生成内容、推荐连接时，用户唯一不可替代的东西可能变成“自己是谁”。如果一个人长期积累了作品、观点、审美和关系，AI 就能更准确地放大他；如果一个人没有主体性，AI 只能生成平均化表达。

\begin{importantbox}{主体性是未来的个人资产}
过去互联网时代，流量和粉丝是资产；AI 社交时代，能被模型理解和代理的主体性会成为资产。它包括作品、表达、价值观、关系和边界。
\end{importantbox}

\subsection{分身的边界}

分身越像用户，风险越高：它可能说错话、代表用户做出不被授权的表达，或者泄露隐私。主体性因此必须包含边界：哪些内容能公开，哪些必须确认，哪些永远不能被分身使用。

\begin{warningbox}{没有边界的分身不是产品，是风险}
如果一个分身只追求“像你”，却没有权限边界、审计和撤回机制，它会变成用户无法控制的代理。AI 社交产品必须同时构建相似性和控制感。
\end{warningbox}

\subsection{本章小结}

主体性是 AI 社交的核心输入。它让分身能代表用户行动，也要求产品提供边界和控制机制。

